\documentclass[10pt,a4paper]{article}
\usepackage[margin=0.6in]{geometry}
\usepackage{amsmath,amssymb,amsfonts,stmaryrd}
\usepackage{xcolor}
\usepackage{enumitem}
\usepackage{tcolorbox}
\usepackage{multicol}
\usepackage{listings}
\usepackage{hyperref}

\definecolor{isabelleblue}{RGB}{0, 70, 140}
\definecolor{isabellekw}{RGB}{120, 0, 120}
\definecolor{boxbg}{RGB}{248, 249, 250}

\tcbset{
  colback=boxbg,
  colframe=isabelleblue,
  fonttitle=\bfseries\large,
  boxrule=0.8pt,
  arc=3pt,
  left=8pt,
  right=8pt,
  top=6pt,
  bottom=6pt
}

\title{\textbf{\huge Defining Recursive Functions on Lists}\\[0.3em]
\large Lecture Notes}
\author{\textbf{Course:} Computer-Assisted Theorem Proving (7CCSACTP)}
\date{}

\begin{document}

\maketitle
\vspace{-1.5em}

\begin{tcolorbox}[title={1. Context \& List Representation Intuition}]
\begin{multicols}{2}
\textbf{Last Week Revision:}
\begin{itemize}[leftmargin=1.2em, itemsep=2pt]
    \item Introduction to Isabelle interface (terms, theorems)
    \item Proof styles: FW, BW, induction, substitution, cases
\end{itemize}

\textbf{This Video Agenda:}
\begin{itemize}[leftmargin=1.2em, itemsep=2pt]
    \item Defining recursive functions on lists (Theory + Isabelle)
\end{itemize}

\columnbreak

\textbf{List Intuition \& Structure:}
\begin{itemize}[leftmargin=1.2em, itemsep=2pt]
    \item Any list $xs$ is constructed via:
    \begin{itemize}[leftmargin=1em, label=$\bullet$]
        \item \texttt{Nil} (\texttt{[]})
        \item \texttt{Cons x xs} ($x \# xs$)
    \end{itemize}
    \item Syntactic sugar:
    \[ a \# [] \approx [a], \qquad [a,b,c] \approx a \# b \# c \# [] \]
    \item Uses: Arrays, Association lists (e.g. $[(\text{Bob}, 23), (\text{Alice}, 34)]$)
    \item \textbf{Finiteness:} Lists in Isabelle are always finite applications of \texttt{Cons}: $x_1 \# x_2 \# \dots \# x_n \# []$.
\end{itemize}
\end{multicols}
\end{tcolorbox}

\vspace{0.3em}

\begin{tcolorbox}[title={2. Recursive Functions on Lists}]
\textbf{Concept:} A function $f$ on lists is recursive if $f(x \# xs)$ is evaluated in terms of $f(xs)$ (calling itself on the tail). A fully-defined recursive function must cover all input patterns (\texttt{[]} and $x \# xs$).

\begin{multicols}{2}
\textbf{Example 1: Sum of List}
\begin{itemize}[leftmargin=1.2em, itemsep=2pt]
    \item \textbf{Signature:} \texttt{sum :: nat list $\Rightarrow$ nat}
    \item \textbf{Equations:}
    \begin{itemize}[leftmargin=1em, label=$\bullet$]
        \item \texttt{sum [] = 0}
        \item \texttt{sum (x \# xs) = x + sum xs}
    \end{itemize}
    \item \textbf{Evaluation trace:}
    \begin{align*}
    \text{sum }[0,1,3] &= 0 + \text{sum }[1,3] \\
    &= 0 + (1 + \text{sum }[3]) \\
    &= 0 + (1 + (3 + \text{sum }[])) \\
    &= 0 + (1 + (3 + 0)) = 4
    \end{align*}
\end{itemize}

\columnbreak

\textbf{Example 2: Append}
\begin{itemize}[leftmargin=1.2em, itemsep=2pt]
    \item \textbf{Signature:} \texttt{append :: 'a list $\Rightarrow$ 'a list $\Rightarrow$ 'a list}
    \item \textbf{Equations:}
    \begin{itemize}[leftmargin=1em, label=$\bullet$]
        \item \texttt{append [] ys = ys}
        \item \texttt{append (x \# xs) ys = x \# (append xs ys)}
    \end{itemize}
    \item \textbf{Evaluation trace:}
    \begin{align*}
    \text{append }[1,2]\;[3] &= 1 \# (\text{append }[2]\;[3]) \\
    &= 1 \# (2 \# (\text{append }[]\;[3])) \\
    &= 1 \# (2 \# [3]) = [1,2,3]
    \end{align*}
\end{itemize}
\end{multicols}
\end{tcolorbox}

\vspace{0.3em}

\begin{tcolorbox}[title={3. Important Rules for Termination \& Simpset}]
\begin{itemize}[leftmargin=1.2em, itemsep=3pt]
    \item \textbf{Termination Requirement:} To guarantee termination, all recursive calls must be on strictly smaller arguments (e.g. $xs$ is structurally smaller than $x \# xs$).
    \item \textbf{Simpset Generation:} Only proven terminating functions will have automatic simplification rules (\texttt{f.simps}) generated and added to Isabelle's simpset.
\end{itemize}
\end{tcolorbox}

\end{document}
